\documentclass[border=8pt]{standalone}
\usepackage{pgfplots}
\pgfplotsset{compat=1.18, trig format plots=deg}
\begin{document}
\begin{tikzpicture}
\begin{axis}[
    width  = 12.2cm, height = 9.4cm,
    view={24}{-58},
    xlabel=$x$, ylabel=$y$, zlabel=$z$,
    xmin=-1.15, xmax=1.15, ymin=-0.9, ymax=0.9, zmin=-0.78, zmax=0.78,
    ztick={-0.6,-0.3,...,0.6},
    title={Part (a): maximizing region
           $E=\{(x,y,z)\mid x^2+2y^2+3z^2\le 1\}$},
    every title/.append style={font=\small},
    axis background/.style={fill=black!3},
]
% ---- the ellipsoid surface ----------------------------------------------
\addplot3[surf, shader=flat, opacity=0.38,
          domain=0:360, samples=37, y domain=0:180, samples y=19,
          color=orange!70!red]
    ({cos(x)*sin(y)}, {sin(x)*sin(y)/sqrt(2)}, {cos(y)/sqrt(3)});

% ---- principal traces of the boundary -----------------------------------
\addplot3[very thick, black, domain=0:360, samples=121, smooth]
    ({cos(x)}, {sin(x)/sqrt(2)}, 0);
\node[font=\scriptsize, black, anchor=north east] at (axis cs:-0.30,-0.62,0)
     {$z{=}0:\ x^2{+}2y^2{=}1$};
\addplot3[thick, black, domain=0:360, samples=121, smooth]
    ({cos(x)}, 0, {sin(x)/sqrt(3)});
\node[font=\scriptsize, black, anchor=north west] at (axis cs:0.75,0.0,-0.60)
     {$y{=}0:\ x^2{+}3z^2{=}1$};
\addplot3[thick, black, domain=0:360, samples=121, smooth]
    (0, {cos(x)/sqrt(2)}, {sin(x)/sqrt(3)});

% ---- semi-axes ------------------------------------------------------------
\addplot3[red,   ultra thick] coordinates {(0,0,0) (1,0,0)}
    node[midway, below, font=\footnotesize, red]{$a{=}1$};
\addplot3[blue,  ultra thick] coordinates {(0,0,0) (0,0.7071,0)}
    node[midway, right, font=\footnotesize, blue]{$b{=}1/\sqrt{2}$};
\addplot3[green!55!black, ultra thick] coordinates {(0,0,0) (0,0,0.5774)}
    node[midway, right, font=\footnotesize, green!55!black]{$c{=}1/\sqrt{3}$};

% ---- vertices -------------------------------------------------------------
\addplot3[only marks, mark=*, mark size=1.6pt, black]
    coordinates {(1,0,0) (-1,0,0) (0,0.7071,0) (0,-0.7071,0) (0,0,0.5774) (0,0,-0.5774)};
\end{axis}
\end{tikzpicture}
\end{document}
